\begin{thebibliography}{99}
\bibitem{ref1} I. A. Alfarisi, A. T. Priandika, and A. S. Puspaningrum, “Penerapan Framework Laravel Pada Sistem Pelayanan Kesehatan (Studi Kasus: Klinik Berkah Medical Center),” J. Ilm. Comput. Sci., vol. 2, no. 1, pp. 1–9, 2023, doi: 10.58602/jics.v2i1.11.
\bibitem{ref2} M. S. Deni Murdiani, “PERBANDINGAN METODOLOGI WATERFALL DAN RAD (RAPID APPLICATION DEVELOPMENT) DALAM PENGEMBANGAN SISTEM INFORMASI,” vol. 4, no. 4, pp. 302–306, 2022.
\bibitem{ref3} H. N. Ningsih, Widya, “PERBANDINGAN MODEL WATERFALL DAN METODE PROTOTYPE UNTUK PENGEMBANGAN APLIKASI PADA SISTEM INFORMASI,” vol. 5, no. 1, pp. 83–95, 2023.
\bibitem{ref4} M. T. P. Perdana, G. Farell, T. Sriwahyuni, and L. Mursyida, “Design and Implementation of a Web-Based Internship Management Information System for Vocational Education,” J. Vokasi Inform., vol. 5, no. 2, pp. 286– 304, 2025.
\bibitem{ref5} M. Nugraha, R. Agus, H. Fathi, and R. Baginda, “Development a Web-Based Student Internship Application Using Laravel Framework,” J. Inf. Technol. Its Util., vol. 6, no. 1, pp. 31–38, 2023, doi: 10.56873/jitu.6.1.5139.
\bibitem{ref6} H. D. Yulianto and R. B. Firdaus, “Perancangan Sistem Informasi Monitoring Magang,” IJIS - Indonesian J. Inf. Syst., vol. 6, no. 2, pp. 130--136, 2021, doi: 10.36549/ijis.v6i2.144.
\bibitem{ref7} F. Moeinzadeh, S. H. R. Ayati, B. Iraj, M. Mortazavi, and V. Vafamehr, “Designing, implementation, and evaluation of internship comprehensive system for assessment and monitoring,” J. Educ. Health Promot., vol. 10, art. no. 93, 2021, doi: 10.4103/jehp.jehp\_626\_20.
\bibitem{ref8} M. Anas, U. M. Putri, and M. Yusuf, “Sistem Informasi Manajemen Praktek Kerja Lapangan ( PKL ) Berbasis Website pada PT . Perkebunan Nusantara IV Regional IV Kota Jambi Menggunakan Extreme Programming,” vol. 4, no. 2, pp. 1059–1069, 2025.
\bibitem{ref9} M. Arkan, R. Razan, I. Arwani, and H. Farisi, “Pengembangan Aplikasi Monitoring dan Evaluasi Kinerja Siswa Magang dari Indonesia Learning Center berbasis Android,” J. Pengemb. Teknol. Inf. dan Ilmu Komput., vol. 7, no. 5, pp. 2091–2100, 2023, [Online]. Available: \url{http://j-ptiik.ub.ac.id}
\bibitem{ref10} A. A. Samudra et al., “PERANCANGAN SISTEM INFORMASI MANAJEMEN PRAKTIK KERJA LAPANGAN ( PKL ) BERBASIS WEB DI SMK NEGERI 1 BUNGO,” vol. 7, no. 5, pp. 3497–3503, 2023.
\bibitem{ref11} B. Bestin, M. G. Addiansyah, R. Sofian, F. Reza, and M. N. Wafa, “Sistem Informasi Monitoring Peserta Magang Berbasis Web Pada Divisi Sales \& Business Development Di Mayar Kota Bandung,” vol. 14, no. 2, pp. 163–170, 2023.
\bibitem{ref12} F. K. Dewi, E. Purwanto, and H. Permatasari, “Sistem Informasi Monitoring Praktik Kerja Lapangan Siswa Berbasis Web Di SMK Ma ’ arif Nu 2 Boyolali,” pp. 637–647, 2024.
\bibitem{ref13} M. Ilmi, D. R. Habibie, and Y. Arifin, “Analisis dan Perancangan Sistem Informasi Monitoring Siswa PKL pada SMK Permata Harapan,” JOINS (Journal Inf. Syst., vol. 8, no. 2, pp. 177–187, 2023, doi: 10.33633/joins.v8i2.9233.
\bibitem{ref14} N. F. Syahidannur Mukminah Pulungan, Reza Febrianti, Titin Lestari, Natra Gurning, “Analisis Teknik Entity-Relationship Diagram Dalam Perancangan Database,” vol. 2, no. 1, pp. 98–102, 2023.
\bibitem{ref15} I. R. Mukhlis and R. Santoso, “Perancangan Basis Data Perpustakaan Universitas Menggunakan MySQL dengan Physical Data Model dan Entity Relationship Diagram,” vol. 4, no. 2, pp. 2–8, 2023, doi: 10.37802/joti.v4i2.330.
\bibitem{ref16} F. Renardi, N. Faizah, and L. Koryanto, “Reimbursement Application Using Rest API Methods and Android-Based Flutter Framework (Case Study: PT. Protonema),” J. Mob. Technol., vol. 1, no. 1, pp. 44–52, 2023, doi: 10.59431/jms.v1i1.133.
\bibitem{ref17} Y. L. Pattinama and I. Susanti, “Implementasi Rest API Web Service Dengan Otentifikasi JSON Web Token Untuk Aplikasi Properti,” vol. 4221, no. April, pp. 81–89, 2023.
\bibitem{ref18} A. S. F. Ferry Kurniawan, Erlin Ayu Khrisnawati, Rizka Hadiwiyanti, “PENGUJIAN SISTEM INFORMASI MANAJEMEN SISWA BERBASIS WEBSITE MENGGUNAKAN METODE BLACK BOX DAN WHITE BOX Ferry,” no. September, pp. 10–11, 2022.
\bibitem{ref19} A. S. Wulandari, A. Saepudin, M. P. Kinanti, Z. Sudesi, A. Saifudin, and Y. Yulianti, “Pengujian Aplikasi Sistem Informasi Akademik Berbasis Web Menggunakan Metode Black Box Testing Equivalence Partitioning,” J. Teknol. Sist. Inf. dan Apl., vol. 5, no. 2, p. 102, 2022, doi: 10.32493/jtsi.v5i2.17561.
\bibitem{ref20} I. Permatasari, F. Adhania, S. A. Putri, and S. R. C. Nursari, “Pengujian Black Box Menggunakan Metode Analisis Nilai Batas pada Aplikasi DANA,” KONSTELASI Konvergensi Teknol. dan Sist. Inf., vol. 3, no. 2, pp. 373–387, 2023, doi: 10.24002/konstelasi.v3i2.8289.
\end{thebibliography}
